% Part: set-theory
% Chapter: z
% Section: arbintersections

\documentclass[../../../include/open-logic-section]{subfiles}

\begin{document}

\olfileid{sth}{z}{arbintersections}
\olsection{পরিশিষ্ট: আবরণ, ধর্মনির্দেশে সেট-গঠন ও ছেদ}

\olref[sth][z][milestone]{sec}-এ আমরা বলেছিলাম,
\olref[sfr][][]{part}-এর স্বতঃস্ফূর্ত আলোচনায় ফিরে
দেখতে যে, সেখানকার কাজগুলি $\Zminus$-এ করা যায় কি না।
কথাটি মেনে থাকলে \emph{একটি} জায়গায় হয়তো অসুবিধা
হয়েছে: ডেডেকিন্ডের \emph{আবরণ}-এর আলোচনায়
\emph{ছেদ}-এর ব্যবহার।

\olref[sfr][infinite][dedekind]{Closure} থেকে মনে করো:
\begin{align*}
  \closureofunder{f}{o} & = \bigcap\Setabs{X}{o \in X \text{ এবং $X$ হলো
$f$-বদ্ধ}}.
\intertext{এর সাধারণ রূপটি নিচের মতো সংজ্ঞা:}
  C & = \bigcap\Setabs{X}{\phi(X)}.
\end{align*}
কিন্তু এখানে সতর্ক হওয়া উচিত। অনানুষ্ঠানিক ধর্মনির্দেশে
সেট-গঠন ব্যর্থ বলে
$\Setabs{X}{\phi(X)}$ সেটটি যে আছে, তার নিশ্চয়তা নেই।
তাহলে এমন সংজ্ঞায় যেন
\emph{অন্যায় সুবিধা} নেওয়া হচ্ছে।

সৌভাগ্যক্রমে, সত্যিই তা নয়; অথবা বর্তমানে লেখা রূপে
যদি তা-ই হয়, তবে সৎ উপায়ে কিছু কাজ করে সংজ্ঞাগুলিকে
বৈধ করা যায়। সেই কাজের ইঙ্গিত
\olref[sth][z][sep]{prop:intersectionsexist}-এ ছিল,
যেখানে ব্যাখ্যা করেছি কেন যেকোনো $A \neq \emptyset$-এর
জন্য $\bigcap A$ আছে। এবার বিষয়টি স্পষ্ট করে লিখি।

বিস্তারগততা ধরে, যদি $C$-কে
$\bigcap\Setabs{X}{\phi(X)}$ হিসেবে সংজ্ঞায়িত করতে চাই,
তবে আমরা আসলে এমন একটি বস্তু $C$ চাই যা নিচের শর্তটি
মেনে চলে:
\begin{equation}\ollabel{bicondelimarbintersection}
	\forall x(x \in C \liff \forall X(\phi(X) \lif x \in X))\tag{*}
\end{equation}
এখন ধরি, \emph{অন্তত একটি} সেট $S$ আছে, যাতে
$\phi(S)$ সত্য। তাহলে
\olref{bicondelimarbintersection} পাওয়ার জন্য
\emph{পৃথকীকরণ} ব্যবহার করে এভাবে $C$ সংজ্ঞায়িত
করলেই চলে:
\[
	C = \Setabs{x \in S}{\forall X(\phi(X) \lif x \in X)}.
\]
এই সংজ্ঞা থেকে যে কাঙ্ক্ষিত
\olref{bicondelimarbintersection} পাওয়া যায়, তা
অনুশীলন হিসেবে যাচাই করো।
%  একটি $x$ নিই। প্রথমে ধরি, নতুন সংজ্ঞা অনুযায়ী $x \in C$; তাহলে $x
%  \in S$ এবং $\forall X(\phi(X) \lif x \in X)$—উভয়ই সত্য; কিন্তু
%  $\phi(S)$ হওয়ায় শুধু $\forall X(\phi(X)
%  \lif x \in X)$ শর্তটিই যথেষ্ট। উল্টোদিকে, ধরি $\forall X(\phi(X) \lif
%  x \in X)$; তাহলে $\phi(S)$ থেকে $x \in S$-ও পাই; অতএব
%  নতুন সংজ্ঞা অনুযায়ী $x \in c$।
ছেদ সংজ্ঞায়িত করতে গিয়ে অনানুষ্ঠানিক ধর্মনির্দেশে
সেট-গঠনের আপাত
ব্যবহার এড়ানোর জন্য এই সাধারণ কৌশল কাজে লাগে।
যে বিশেষ ক্ষেত্রে আলোচনা শুরু হয়েছিল, অর্থাৎ
$\closureofunder{f}{o}$-এর ক্ষেত্রে, কৌশলটি এভাবে
প্রয়োগ করা যায়।
\olref[sfr][infinite][dedekind]{closureproperties}-এর
প্রমাণের শুরুতে আমরা দেখেছিলাম যে,
$o \in \ran{f}\cup\{o\}$ এবং $\ran{f} \cup \{o\}$
$f$-বদ্ধ। তাই চাওয়া বস্তুটি এভাবে
সংজ্ঞায়িত করতে পারি:
\[
	\closureofunder{f}{o} = \Setabs{x \in \ran{f} \cup \{o\}}{(\forall X \ni o)(X \text{ হলো $f$-বদ্ধ} \lif x \in X)}.
\]

\end{document}
